% !TEX root = ../main.tex

\chapter{参考值的计算与可靠性说明}\label{ch:refval}
  本附录说明第\ref{ch:numerical}章中各算例参考值的计算方法.
  对具有显式解的方程, 直接使用解析公式.
  对没有简单闭式解的方程, 使用与正文 RFM 求解器无关的数值方法计算参考值,
  并给出相应的置信区间或误差估计.

  \section{Heat 方程的显式参考值}\label{sec:ref_heat}
    Heat 方程的终端条件取为
    \begin{equation}
      g(x)=\exp\!\left(-\frac{\lVert x\rVert^2}{d}\right),
      \qquad x_0=0,\qquad T=1.
      \label{eq:heatterm}
    \end{equation}
    更一般地, 若终端条件写为
    \begin{equation*}
      g(x)=\exp\!\left(-\frac{\lVert x\rVert^2}{a}\right),
    \end{equation*}
    则由 Feynman--Kac 公式可得
    \begin{equation}
      u(t,x)
      =
      \left(\frac{a}{a+2(T-t)}\right)^{d/2}
      \exp\!\left(
              -\frac{\lVert x\rVert^2}{a+2(T-t)}
      \right).
      \label{eq:numerical_heat_solution}
    \end{equation}
    因而在 $x_0=0$, $a=d$ 时,
    \begin{equation}
      Y_0 = \left(\frac{d}{d+2}\right)^{d/2}.
      \label{eq:heaty0}
    \end{equation}
    代入不同维度得到表~\ref{tab:ref_heat}.
    这些值由显式公式直接计算, 不包含额外的数值离散误差.

    \begin{table}
      \centering
      \caption{Heat 方程参考值}
      \label{tab:ref_heat}
      \begin{tabular}{cc}
        \toprule
        $d$ & $Y_0$ \\
        \midrule
        20 & 0.385543289429532 \\
        30 & 0.379812405815246 \\
        50 & 0.375116802253964 \\
        100 & 0.371527882126961 \\
        200 & 0.369711212329119 \\
        300 & 0.369102310996184 \\
        \bottomrule
      \end{tabular}
    \end{table}

  \section{HJB-LQ 方程的积分参考值}\label{sec:ref_hjb}
    HJB-LQ 方程的终端条件取为
    \begin{equation}
      g(x) = \log \left(\frac{1+\lVert x \rVert^{2}}{2}\right).
      \label{eq:hjbterm}
    \end{equation}
    HJB-LQ 方程对应的 BSDE 为
    \begin{equation}
      \dif Y_t=\frac{1}{2}\lVert Z_t\rVert^2\dif t+Z_t^{\operatorname{T}}\dif W_t.
      \label{eq:numerical_hjblq_bsde}
    \end{equation}
    令 $U_t=\exp(-Y_t)$, 则 $U_t$ 为鞅.
    于是
    \begin{equation}
      Y_0
      =
      -\log
      \mathbb{E}
      \left[
        \frac{2}{1+2T\chi_d^2}
      \right].
      \label{eq:hjby0}
    \end{equation}
    因此参考值可以通过下面的一维积分计算:
    \begin{equation}
      \mathbb{E}
      \left[
        \frac{2}{1+2T\chi_d^2}
      \right]
      =
      \int_{0}^{\infty}
      \frac{2}{1+2Ty} f_{\chi_d^2}(y)\,\dif y.
      \label{eq:hjbint}
    \end{equation}
    本文采用自适应 Simpson 积分在有限区间上计算该积分.
    当积分上限从 $300$ 增加到 $500$ 时, 结果变化不超过 $5\times 10^{-12}$,
    因而在本文所列精度下, 截断误差可以忽略.
    最终结果见表~\ref{tab:ref_hjb}.

    \begin{table}
      \centering
      \caption{HJB-LQ 方程参考值}
      \label{tab:ref_hjb}
      \begin{tabular}{cc}
        \toprule
        $d$ & $Y_0$ \\
        \midrule
        20 & 2.921014735738689 \\
        30 & 3.351223237182463 \\
        50 & 3.882006682195226 \\
        100 & 4.590161724604434 \\
        200 & 5.290814736008237 \\
        300 & 5.698781231260821 \\
        \bottomrule
      \end{tabular}
    \end{table}

  \section{Black--Scholes 方程的 Monte Carlo 参考值}\label{sec:ref_bsm}
    对 Black--Scholes 方程, 终端收益函数为高维算术平均型欧式看涨期权
    \begin{equation}
      g(x)=\max\!\left(\frac{1}{d}\sum_{m=1}^{d}x_m-K,0\right).
      \label{eq:bsterm}
    \end{equation}
    经典一维欧式期权定价已经有成熟的 Black--Scholes 基本方程和闭式公式\cite{scholesPricingOptionsCorporate}.
    但这里的算术平均收益依赖多个标的资产, 因而该问题没有类似一维 Black--Scholes 公式的简单闭式解.
    因此本文采用独立的大样本 Monte Carlo 方法估计其初值,
    并使用几何平均期权的闭式定价公式作为控制变量.

    记算术平均收益对应的贴现随机变量为 $A$,
    几何平均收益对应的贴现随机变量为 $G$.
    由于几何平均
    \begin{equation}
      \bar{G}_T = \exp\!\left(\frac{1}{d}\sum_{m=1}^{d}\log X_T^{m}\right),
      \label{eq:bsgbar}
    \end{equation}
    仍服从对数正态分布, 因而 $\mathbb{E}[G]$ 可以显式计算.
    于是采用控制变量估计量
    \begin{equation}
      A^{\ast}=A-\beta\bigl(G-\mathbb{E}[G]\bigr),
      \label{eq:bsdsigma}
    \end{equation}
    同时结合对偶路径进一步降低方差.

    本文使用 $10^6$ 条路径计算参考值, 结果和 $95\%$ 置信区间见表~\ref{tab:ref_bsm}.

    \begin{table}
      \centering
      \caption{Black--Scholes 方程 Monte Carlo 参考值}
      \label{tab:ref_bsm}
      \begin{tabular}{ccc}
        \toprule
        $d$ & $Y_0$ & $95\%$ 置信区间 \\
        \midrule
        20 & 0.05173661 & $[0.05172011,\ 0.05175310]$ \\
        30 & 0.05021835 & $[0.05020434,\ 0.05023235]$ \\
        50 & 0.04921751 & $[0.04920643,\ 0.04922860]$ \\
        100 & 0.04881270 & $[0.04880488,\ 0.04882053]$ \\
        200 & 0.04877301 & $[0.04876747,\ 0.04877854]$ \\
        300 & 0.04876993 & $[0.04876541,\ 0.04877445]$ \\
        \bottomrule
      \end{tabular}
    \end{table}

  \section{Allen--Cahn 方程的有限差分参考值}\label{sec:ref_allencahn}
    对 Allen--Cahn 方程, 由于其终端条件和扩散过程都具有径向对称性,
    该问题可以转化为径向变量上的一维半线性抛物方程.
    记 $r=\lVert x\rVert$, 则解可写为 $u(t,x)=v(t,r)$.
    在反向时间变量 $\tau=T-t$ 下, 参考解满足
    \begin{equation}
      \frac{\partial v}{\partial \tau}
      =
      \frac{\partial^2 v}{\partial r^2}
      + \frac{d-1}{r}\frac{\partial v}{\partial r}
      + v - v^3,
      \qquad
      v(0,r)=\frac{1}{2+0.4r^2}.
      \label{eq:acv}
    \end{equation}
    Allen--Cahn 方程是高维半线性 PDE 算法中常用的反应--扩散型测试问题\cite{hanSolvingHighdimensionalPartial2018}.
    本文在区间 $r\in[0,20]$ 上采用径向有限差分离散,
    对扩散项作隐式处理, 对非线性反应项作显式处理.
    在 $r=0$ 处施加径向对称边界条件, 在远端边界取零值近似.

    为评估离散误差, 分别使用两组网格
    \begin{equation*}
      (\Delta r,\Delta \tau)=(0.025,\ 5\times 10^{-5}),
    \end{equation*}
    \begin{equation*}
      (\Delta r,\Delta \tau)=(0.0125,\ 2.5\times 10^{-5})
    \end{equation*}
    进行计算, 并将两组结果的差作为误差估计.
    最终采用细网格结果作为参考值, 结果见表~\ref{tab:ref_allencahn}.

    \begin{table}
      \centering
      \caption{Allen--Cahn 方程有限差分参考值}
      \label{tab:ref_allencahn}
      \begin{tabular}{ccc}
        \toprule
        $d$ & $Y_0$ & 网格差估计 \\
        \midrule
        20 & 0.20635880 & $2.08\times 10^{-6}$ \\
        30 & 0.15199385 & $1.59\times 10^{-6}$ \\
        50 & 0.09908593 & $3.19\times 10^{-6}$ \\
        100 & 0.05278464 & $2.81\times 10^{-6}$ \\
        200 & 0.02724542 & $1.79\times 10^{-6}$ \\
        300 & 0.01835709 & $1.29\times 10^{-6}$ \\
        \bottomrule
      \end{tabular}
    \end{table}

    此外, 当边界截断半径从 $15$ 增加到 $30$ 时结果几乎不变,
    因而边界截断误差小于当前网格误差.
